% Part: lambda-calculus
% Chapter: syntax
% Section: alpha

\documentclass[../../../include/open-logic-section]{subfiles}

\begin{document}

\olfileid{lam}{syn}{alp}
\olsection{Konversi-$\alpha$}

Apa sesambungane $\lambd[x][x]$ lan
$\lambd[y][y]$? Loro-lorone makili
fungsi identitas. Mesthi wae, minangka
suku, sintaksise beda. Bedane mung
jeneng variabel kaiket; suku siji bisa
dipikolehi kanthi ngganti jeneng
variabel kaiket ing suku sijine.
Iki diarani \emph{konversi-$\alpha$}.

\begin{defn}[Owah-owahan variabel kaiket, $\aconvone$]
  Yen suku $M$ ngemot muncule
  $\lambd[x][N]$, $x \neq y$,
  $y \notin \FV{N}$, lan
  $\Subst{N}{y}{x}$ katetepake,
  mula ngganti muncule iku nganggo
  \begin{equation*}
    \lambd[y][\Subst{N}{y}{x}]
  \end{equation*}
  nganti ngasilake $M'$ diarani
  \emph{owah-owahan variabel kaiket},
  ditulis $M \redone[\alpha] M'$.
\end{defn}

\begin{defn}[Kompatibilitas relasi]
  Relasi $R$ ing suku diarani
  \emph{kompatibel} yen nyukupi
  syarat-syarat iki:
  \begin{enumerate}
  \item Yen $R N N'$, mula
    $R \lambd[x][N] \lambd[x][N']$.
  \item Yen $R P P'$, mula $R (PQ) (P'Q)$.
  \item Yen $R Q Q'$, mula $R (PQ) (PQ')$.
  \end{enumerate}
\end{defn}

Dadi, definisi kasebut bisa
dinyatakake maneh:
\begin{defn}[Owah-owahan variabel kaiket, $\aconvone$]
  \emph{Owah-owahan variabel kaiket}
  ($\redone[\alpha]$) yaiku relasi
  kompatibel paling cilik ing suku
  kang nyukupi syarat iki:
  \begin{align*}
    &\lambd[x][N] \redone[\alpha] \lambd[y][\Subst{N}{y}{x}] && \text{yen
      $x \neq y$, $y \notin \FV{N}$} \\
    & &&\text{lan $\Subst{N}{y}{x}$ katetepake}
  \end{align*}
\end{defn}

Tembung ``paling cilik'' tegese relasi
iku mung ngemot pasangan kang diwajibake
dening kompatibilitas lan syarat tambahan
ing ndhuwur. Mula relasi iki uga bisa
ditetepake kaya mangkene:

\begin{defn}[Owah-owahan variabel kaiket, $\aconvone$] \ollabel{defn:aconvone}
  \emph{Owah-owahan variabel kaiket}
  ($\aconvone$) ditetepake kanthi
  induktif kaya mangkene:
  \begin{enumerate}
  \item Yen $N \aconvone N'$, mula
    $\lambd[x][N] \aconvone
    \lambd[x][N']$. \ollabel{defn:aconvone1}
  \item Yen $P \aconvone P'$, mula
    $(PQ) \aconvone (P'Q)$. \ollabel{defn:aconvone2}
  \item Yen $Q \aconvone Q'$, mula
    $(PQ) \aconvone (PQ')$. \ollabel{defn:aconvone3}
  \item Yen $x \neq y$, $y \notin \FV{N}$,
    lan $\Subst{N}{y}{x}$ katetepake,
    mula $\lambd[x][N] \redone[\alpha]
    \lambd[y][\Subst{N}{y}{x}]$.
    \ollabel{defn:aconvone4}
  \end{enumerate}
\end{defn}

Definisi-definisi iki ekuivalen,
nanging buktine dadi latihan.
Sabanjure kita nggunakake
definisi induktif.

\begin{defn}[Konversi-$\alpha$, $\aconv$]
  \emph{Konversi-$\alpha$} ($\aconv$)
  yaiku relasi refleksif lan transitif
  paling cilik ing suku kang ngemot
  $\aconvone$.
\end{defn}

Kaya ing ndhuwur, ``paling cilik''
ateges relasi iku mung ngemot pasangan
kang diwajibake dening refleksivitas,
transitivitas, lan $\aconvone$.
Mula kita oleh definisi ekuivalen iki:

\begin{defn}[Konversi-$\alpha$, $\aconv$] \ollabel{defn:aconv}
  \emph{Konversi-$\alpha$} ($\aconv$)
  ditetepake kanthi induktif kaya mangkene:
  \begin{enumerate}
  \item Yen $P \aconv Q$ lan $Q \aconv R$,
    mula $P \aconv R$.
    \ollabel{defn:aconv1}
  \item Yen $P \aconvone Q$, mula
    $P \aconv Q$. \ollabel{defn:aconv2}
  \item $P \aconv P$. \ollabel{defn:aconv3}
  \end{enumerate}
\end{defn}

\begin{ex}
  $\lambd[x][f x]$ bisa dikonversi-$\alpha$
  dadi $\lambd[y][f y]$, lan uga
  kosok baline. Miturut intuisi,
  loro-lorone fungsi kang nampa
  argumen lan mbalekake asil $f$
  ing argumen iku, kanthi $f$
  ngrujuk marang variabel lingkungan.

  $\lambd[x][f x]$ ora bisa
  dikonversi-$\alpha$ dadi
  $\lambd[x][g x]$. Miturut intuisi,
  suku-suku iku ngrujuk marang
  variabel lingkungan $f$ lan~$g$
  kanthi urut. Mula fungsine beda:
  tumindake beda ing lingkungan
  kang biji $f$ lan $g$ beda.
\end{ex}

\begin{prob}
  Apa pasangan suku ing ngisor iki
  bisa dikonversi-$\alpha$?
  \begin{enumerate}
  \item $\lambd[x][\lambd[y][x]]$ lan $\lambd[y][\lambd[x][y]]$
  \item $\lambd[x][\lambd[y][x]]$ lan $\lambd[c][\lambd[b][a]]$
  \item $\lambd[x][\lambd[y][x]]$ lan $\lambd[c][\lambd[b][a]]$
  \end{enumerate}
\end{prob}

\begin{lem}\ollabel{lem:fv-one}
  Yen $P \aconvone Q$, mula
  $\FV{P} = \FV{Q}$.
\end{lem}

\begin{proof}
  Kanthi induksi tumrap !!{derivation}
  saka $P \aconvone Q$.
  \begin{enumerate}
  \item Yen aturan pungkasan yaiku
    \olref{defn:aconvone4}, mula
    $P$ awujud $\lambd[x][N]$
    lan $Q$ awujud
    $\lambd[y][\Subst{N}{y}{x}]$,
    kanthi $x \neq y$,
    $y \notin \FV{N}$ lan
    $\Subst{N}{y}{x}$ katetepake.
    Kita mbedakake kasus manut
    apa $x \in \FV{N}$:
    \begin{enumerate}
    \item Yen $x \in \FV{N}$, mula:
      \begin{align*}
        \FV{\lambd[y][\Subst{N}{y}{x}]} &=
        \FV{\Subst{N}{y}{x}} \setminus \{y\} \\
        & = ((\FV{N} \setminus \{x\}) \cup \{y\}) \setminus \{y\}
         && \text{miturut \olref[sub]{thm:infv}} \\
        & = \FV{N} \setminus \{x\} \\
        & = \FV{\lambd[x][N]}
      \end{align*}
    \item Yen $x \notin \FV{N}$, mula:
      \begin{align*}
        \FV{\lambd[y][\Subst{N}{y}{x}]}
        & = \FV{\Subst{N}{y}{x}} \setminus \{y\} \\
        & = \FV{N} \setminus \{y\}
         && \text{miturut \olref[sub]{thm:notinfv}} \\
        & = \FV{N} \setminus \{x\} \\
        & = \FV{\lambd[x][N]}.
      \end{align*}
    \end{enumerate}
  \item Telung kasus liyane dadi latihan.
  \end{enumerate}
\end{proof}

\begin{prob}
  Rampungna bukti \olref[lam][syn][alp]{lem:fv-one}.
\end{prob}

\begin{lem}\ollabel{lem:inv}
  Yen $P \aconvone Q$, mula
  $Q \aconvone P$.
\end{lem}

\begin{proof}
  Kanthi induksi tumrap !!{derivation}
  saka $P \aconvone Q$.
  \begin{enumerate}
  \item Yen aturan pungkasan yaiku
    \olref{defn:aconvone4}, mula
    $P$ awujud $\lambd[x][N]$
    lan $Q$ awujud
    $\lambd[y][\Subst{N}{y}{x}]$,
    kanthi $x \neq y$,
    $y \notin \FV{N}$ lan
    $\Subst{N}{y}{x}$ katetepake.
    Kapisan, $x \notin
    \FV{\Subst{N}{y}{x}}$ miturut
    \olref[sub]{thm:clr}, amarga
    $x \neq y$. Miturut
    \olref[sub]{thm:inv},
    $\Subst{\Subst{N}{y}{x}}{x}{y}$
    katetepake lan padha karo~$N$.
    Dadi, miturut
    \olref{defn:aconvone4},
    $\lambd[y][\Subst{N}{y}{x}]
    \aconvone
    \lambd[x][\Subst{\Subst{N}{y}{x}}{x}{y}]
    = \lambd[x][N]$.
  \item Telung kasus kompatibilitas
    ngetutake hipotesis induksi.
  \end{enumerate}
\end{proof}

\begin{prob}
  Rampungna bukti \olref[lam][syn][alp]{lem:inv}.
\end{prob}

\begin{thm}
  Konversi-$\alpha$ iku relasi
  ekuivalensi ing suku, yaiku
  refleksif, simetris, lan transitif.
\end{thm}

\begin{proof}
  \begin{enumerate}
  \item Kanggo saben suku $M$,
    $M$ bisa diowahi dadi $M$
    kanthi \emph{nol} owah-owahan
    variabel kaiket.
  \item Yen $P$ dikonversi-$\alpha$
    dadi $Q$ kanthi runtutan owah-owahan
    variabel kaiket, saka $Q$ kita
    bisa mbalikke owah-owahan mau
    miturut \olref{lem:inv}
    kanthi urutan kosok balen kanggo
    ngasilake~$P$.
  \item Yen $P$ dikonversi-$\alpha$
    dadi $Q$ kanthi sawijining runtutan,
    lan $Q$ dadi $R$ kanthi runtutan liyane,
    kita bisa ngowahi $P$ dadi $R$
    kanthi nindakake runtutan kapisan
    banjur runtutan kapindho.
  \end{enumerate}
\end{proof}

Sabanjure, $M$ lan $N$ diarani
\emph{ekuivalen-$\alpha$}, $M
\aeq N$, yen lan mung yen $M$
bisa dikonversi-$\alpha$ dadi $N$.
Miturut asil ing ndhuwur, iki
ekuivalen karo $N$ bisa
dikonversi-$\alpha$ dadi~$M$.

\begin{thm}\ollabel{thm:fv}
  Yen $M \aeq N$, mula
  $\FV{M} = \FV{N}$.
\end{thm}

\begin{proof}
  Langsung saka \olref{lem:fv-one}
  lan definisi tutupan refleksif-transitif.
\end{proof}

\begin{lem}\ollabel{lem:sub:R}
  Yen $R \aeq R'$ lan
  $\Subst{M}{R}{y}$ katetepake,
  mula $\Subst{M}{R'}{y}$
  katetepake lan ekuivalen-$\alpha$
  karo $\Subst{M}{R}{y}$.
\end{lem}

\begin{proof}
  Latihan.
\end{proof}

\begin{prob}
  Buktikna \olref[lam][syn][alp]{lem:sub:R}.
\end{prob}

Elinga yen ing \olref[sub]{sec},
substitusi ora katetepake ing sawetara
kasus. Nanging kanthi konversi-$\alpha$
ing suku, kita bisa ngganti jeneng
variabel kaiket supaya substitusi
katetepake. Asile njaga
ekuivalensi-$\alpha$, kaya kang
dituduhake teorema iki:

\begin{thm}\ollabel{thm:sub}
  Kanggo sembarang $M$, $R$, lan~$y$,
  ana $M'$ saengga $M \aeq M'$
  lan $\Subst{M'}{R}{y}$ katetepake.
  Saliyane iku, yen ana pasangan liya
  $M'' \aeq M$ lan $R''$ kang
  $\Subst{M''}{R''}{y}$ katetepake
  sarta $R'' \aeq R$, mula
  $\Subst{M'}{R}{y} \aeq
  \Subst{M''}{R''}{y}$.
\end{thm}

\begin{proof}
  Kanthi induksi tumrap pambentukan $M$:
  \begin{enumerate}
  \item $M$ iku variabel~$z$: latihan.
  \item Umpamakna $M$ awujud
    $\lambd[x][N]$. Pilih variabel
    $z$ kang beda saka $x$ lan $y$,
    kanthi $z \notin \FV{N}$
    lan $z \notin \FV{R}$.
    Miturut hipotesis induksi,
    ana $N'$ saengga $N' \aeq N$
    lan $\Subst{N'}{z}{x}$ katetepake.
    Mula $\lambd[x][N] \aeq
    \lambd[x][N']$ uga, miturut
    \olref{defn:aconvone}\olref{defn:aconvone1}.
    Saiki $\lambd[x][N']
    \aeq \lambd[z][\Subst{N'}{z}{x}]$
    miturut
    \olref{defn:aconvone}\olref{defn:aconvone4}.
    Iki bisa ditindakake amarga
    $z \ne x$, $z \notin \FV{N'}$,
    lan $\Subst{N'}{z}{x}$
    katetepake. Pungkasan,
    $\Subst{\lambd[z][\Subst{N'}{z}{x}]}{R}{y}$
    mbutuhake $z \neq y$
    lan $z \notin \FV{R}$.
    Syarat yen substitusi ing awak
    abstraksi uga katetepake
    durung dituduhake ing langkah iki.

    Saliyane iku, yen ana $N''$
    lan $R''$ liyane kang nyukupi
    syarat-syarat padha,
    \begin{multline*}
      \Subst{(\lambd[z][\Subst{N''}{z}{x}])}{R''}{y} \\
    \begin{aligned}
      &= \lambd[z][\Subst{\Subst{N''}{z}{x}}{R''}{y}] \\
      &\aeq \lambd[z][\Subst{\Subst{N''}{z}{x}}{R}{y}] && \text{miturut
        \olref{lem:sub:R}}\\
      &\aeq\lambd[z][\Subst{\Subst{N'}{z}{x}}{R}{y}]
      && \text{miturut hipotesis
        induksi}\\
      &=\Subst{(\lambd[z][\Subst{N'}{z}{x}])}{R}{y}
    \end{aligned}
    \end{multline*}
    \item $M$ awujud $(PQ)$: latihan.
  \end{enumerate}
\end{proof}

\begin{prob}
  Rampungna bukti \olref[lam][syn][alp]{thm:sub}.
\end{prob}

\begin{cor}\ollabel{cor:sub}
  Kanggo sembarang $M$, $R$, lan~$y$,
  ana pasangan $M'$ lan $R'$
  saengga $M' \aeq M$, $R \aeq R'$,
  lan $\Subst{M'}{R'}{y}$
  katetepake. Yen ana pasangan liya
  $M'' \aeq M$ lan $R'' \aeq R'$
  kanthi $\Subst{M''}{R''}{y}$
  katetepake, mula
  $\Subst{M'}{R'}{y} \aeq
  \Subst{M''}{R''}{y}$.
\end{cor}

\begin{proof}
  Langsung saka \olref{thm:sub}.
\end{proof}

\end{document}

